\subsection{Principles of omniscience}
The so-called principles of omniscience are all weaker than the law of excluded middle (LEM), and help measure how close a logical system is to satisfying LEM \cite{HannesDiener, ReverseMathsBishop}.
In this section, we will show that two such principles hold (MP and LLPO), 
and that another one fails (WLPO).

\begin{theorem}[The negation of the weak lesser principle of omniscience ($\neg$WLPO)]\label{NotWLPO}
  \[
    \neg \forall_{\alpha:2^\N} 
    ((\forall_{n:\N}\, \alpha_n = 0 ) \vee \neg (\forall_{n:\N}\, \alpha_n = 0))
  \]
\end{theorem}
\begin{proof}
  We will prove that any decidable property of binary sequences is determined by a finite prefix of fixed length, contradicting $\forall_{n:\N}\, \alpha_n=0$ being decidable for all $\alpha$. 
  Indeed assume $f:2^\mathbb N \to 2$ such that 
  $f(\alpha) = 0$ if and only if $\forall_{n:\mathbb N}\ \alpha_n= 0$. 
  By \Cref{AxStoneDuality}, there is some $c:2[\N]$ with 
  $f(\alpha) = 0$ if and only if $\alpha(c) = 0$. 
  There exists $k:\N$ such that $c$ is expressed in terms of the generators $(g_n)_{n\leq k}$. 
  Now consider $\beta,\gamma:2^\N$ given by 
  $\beta(g_n) = 0$ for all $n:\mathbb N$ and
  $\gamma(g_n) = 0$ if and only if $n\leq k$. 
  As $\beta$ and $\gamma$ are equal on $(g_n)_{n\leq k}$, we have $\beta(c) = \gamma(c)$. 
  However, $f(\beta) = 0$ and $f(\gamma) = 1$, giving a contradiction. 
\end{proof}

\begin{theorem}
  For all $\alpha:\Noo$, we have that 
  \[
    (\neg (\forall_{n:\mathbb N}\, \alpha_n= 0)) \to \Sigma_{n:\mathbb N}\, \alpha_n= 1
  \]
\end{theorem}
\begin{proof}
  By \Cref{ClosedPropAsSpectrum}, we have that $\neg(\forall_{n:\N}\, \alpha_n = 0)$ implies that 
  $\Sp(2/(\alpha_n)_{n:\N})$ is empty. 
  Hence $2/(\alpha_n)_{n:\N}$ is trivial by \Cref{SpectrumEmptyIff01Equal}. 
  Then there exists $k:\N$ such that $\bigvee_{i\leq k} \alpha_i = 1$. 
  As $\alpha_i = 1$ for at most one $i:\N$, 
  there exists a unique $n:\mathbb N$ with $\alpha_n = 1$. 
\end{proof}

\begin{corollary}[Markov's principle (MP)]\label{MarkovPrinciple}
  For all $\alpha:2^\mathbb N$, we have that 
  \[
    (\neg (\forall_{n:\mathbb N}\, \alpha_n= 0)) \to \Sigma_{n:\mathbb N}\, \alpha_n= 1
  \]
\end{corollary}
\begin{proof}
  Given $\alpha:2^\mathbb N$, consider the sequence $\alpha':\Noo$ satisfying $\alpha'_n = 1$ if and only if
  $n$ is minimal with $\alpha_n = 1$. Then apply the above theorem.
\end{proof}

\begin{theorem}[The lesser limited principle of omniscience (LLPO)]\label{LLPO}
  For all $\alpha:\N_\infty$, 
  we have that
  \[\label{eqnLLPO}
    (\forall_{k:\N}\, \alpha_{2k} = 0)  \vee (\forall_{k:\N}\, \alpha_{2k+1} = 0)
  \]
\end{theorem}
\begin{proof}
  Define $f:B_\infty \to B_\infty \times B_\infty$ on generators as follows
  \[\label{eqnLLPOProofMap}
    f(g_n) =\begin{cases}
      (g_k,0) \text{ if } n = 2k\\
      (0,g_k) \text{ if } n = 2k+1\\
    \end{cases}
  \]
  Note that $f$ is a well-defined morphism in $\Boole$ as 
  $f(g_n) \wedge f(g_m) = 0$ whenever $m\neq n$. 
  We claim that $f$ is injective. 
  If $I\subseteq \N$, write 
  $ I_0 =\{k\ |\ 2k \in I\}, 
    I_1 =\{k\ |\ 2k+1 \in I\}
  $.
  Recall that any $x:B_\infty$ is of the form 
  $\bigvee_{i\in I} g_i$ or $\bigwedge_{i\in I} \neg g_i$ for some finite set $I$. 
  \begin{itemize}
    \item If $x = \bigvee_{i\in I} g_i$, then 
      $f(x) = (\bigvee_{i\in I_0}g_{i}, \bigvee_{i\in I_1}g_i)$. 
      So if $f(x) = 0$, then $I_0=I_1 = I = \emptyset$ and $x = 0$. 
    \item Suppose 
      $x = \bigwedge_{i\in I} \neg g_i$.
      Then $f(x) = (\bigwedge_{i\in I_0} \neg g_i, \bigwedge_{i\in I_1} \neg g_i)$, 
      so $f(x) \neq 0$. 
  \end{itemize}
  By \Cref{SurjectionsAreFormalSurjections},
  we have that $f$ corresponds to a surjection 
  $s:\Noo + \Noo \to \Noo$.
  Thus for $\alpha : \Noo$, 
  there exists some $x:\Noo + \Noo$ such that $s(x) = \alpha$. 
  If $x = \mathrm{inl}(\beta)$, 
  then for any $k:\N$ we have that 
\[\alpha_{2k+1} = s(x)_{2k+1} = x(f(g_{2k+1})) = \mathrm{inl}(\beta) (0,g_k)  = \beta(0) = 0.\]
  Similarly, if $x = \mathrm{inr}(\beta)$, we have that $\alpha_{2k} = 0$ for all $k:\N$. 
\end{proof}
The surjection $s:\Noo + \Noo \to \Noo$ above does not have a section. Indeed:
\begin{lemma}
  The function $f$ defined above does not have a retraction. 
\end{lemma}
\begin{proof}
  Suppose $r:B_\infty \times B_\infty \to B_\infty$ is a retraction of $f$. 
  Then $r(0,g_k) = g_{2k+1}$ and $r(g_k,0) = g_{2k}$. 
  %Note that $r(0,1):B_\infty$ is expressible using only finitely many generators $(g_n)_{n\leq N}$.
  Note that $r(0,1) \geq r(0,g_k) = g_{2k+1}$ for all $k:\N$. 
  As a consequence, $r(0,1)$ is of the form $\bigwedge_{i\in I} \neg g_i$ for some finite set $I$.
  %, and by \Cref{BinftyTermsWriting}, 
  %$r(0,1)$ corresponds to a cofinite subset of $\N$. % = \bigwedge_{i:I_0} \neg p_i$, where $i\leq N$ for $i\in I_0$. 
  By similar reasoning so is $r(1,0)$. % corresponds to a cofinite subset of $\N$. 
%  But the intersection of cofinite subsets is cofinite, while 
  But then %this contradicts
  \[r(0,1) \wedge r(1,0) = r( (1,0) \wedge (0,1)) = r(0,0) = 0.\]
  This is a contradiction.
\end{proof}


We finish with an equivalent formulation of LLPO:


\begin{lemma}\label{corAlternativeLLPO}
  Let $(\phi_n)_{n:\N}, (\psi_m)_{m:\N}$ be families of decidable propositions indexed over $\N$.
  We then have 
  \begin{equation}
    (\forall_{n:\N} \forall_{m:\N} (\phi_n \vee \psi_m) )
    \leftrightarrow
    ((\forall_{n:\N} \phi_n) \vee (\forall_{m:\N} \psi_m) )
  \end{equation}
\end{lemma}
\begin{proof}
  See \cite{HannesDiener, ReverseMathsBishop}
\end{proof}
\begin{proof}
  Note that the implication from right to left in the above equation always holds.
  Assume that for all $m,n:\mathbb N$ we have $\phi_n\vee \psi_m$ 
  Consider the sequence $\alpha:2^\mathbb N$ where $\alpha(2n) = 0$ iff $\phi_n$ and 
  $\alpha(2m+1) = 0$ iff $\psi_m$. 
  Let $\beta:\Noo$ be such that $\beta(i) = 1$ iff $i$ is minimal with $\alpha(i) = 1$
  By LLPO, we have that 
  $\beta$ is $0$ on all odd entries or on all even entries. 
  Suppose that $\beta$ hits $0$ on all odd entries. 
  We will show $\psi_m$ for all $m:\N$. 
  As $\beta(2m+1) = 0$, there are two options:
  \begin{itemize}
    	\item If $\alpha(l)=0$ for all $l\leq 2m+1$. Then in particular $\alpha(2m+1)=0$ and $\psi_m$ holds.
	\item Otherwise there is some $l<2m+1$ with $\beta(l) = 1$. 
  As $\beta$ hits $0$ on odd entries, $l$ is even. 
  So $\alpha(2n) = 1$ for $n = \frac{l}2$, meaning that $\neg \phi_n$. 
  By assumption, $\phi_n \vee \psi_m$ holds, hence $\psi_m$ must hold. 
  Thus for all $m:\N$, we have $\psi_m$ if $\beta$ hits $0$ on all odd entries. 
  By a symmetric argument, if $\beta$ hits $0$ on all even entries, we have $\phi_n$ for all $n:\N$. 
  We conclude that 
  $((\forall_{n:\N} \phi_n) \vee (\forall_{m:\N} \psi_m) )$ 
  as required. 
  \end{itemize}
\end{proof}

\begin{remark}
Note that the above statement implies LLPO as $\alpha(2n) =0 \vee \alpha(2m+1) =0$ for all $n,m:\mathbb N$ if $\alpha:\Noo$. 
\end{remark}
